\chapter*{Conclusion}
\addcontentsline{toc}{chapter}{Conclusion}
\markboth{Conclusion}{Conclusion}

Three questions drove this booklet, and three answers emerged.

\textbf{What is data?}
Data is not a fixed substance but a relational one: a fact, signal, or
symbol that is counted as data only within a framework that makes it
relevant. The same string of characters --- a page of Marx --- yields
entirely different data to a child, a historian, and a political
economist, because each brings a different framework to bear. Data
without a framework is not yet data; it is noise.

\textbf{How does something come to stand for something else?}
Through representation: the brain's economy of storing an index rather
than the full event, a word rather than the object, a name rather than
the history it points to. The four criteria --- connection to the real
world, a causal role in behaviour, exclusivity of that role, and a
teleological standard for correctness --- give us a way to ask, for
any candidate, whether it genuinely represents or merely accompanies.

\textbf{What is still open?}
How an abstract representation, once it exists, comes to be carried by
a physical signal. That question opens Chapter~3, which is still being
written. The full answer --- voltages, encoding schemes, and the path
from an abstract bit to a measurable physical state --- is the subject
of the volumes that follow.

\medskip

The project repository (\url{https://github.com/papyrxis/arliz})
tracks progress on the remaining chapters and booklets.
